\chapter*{Resumen Ejecutivo}
\addcontentsline{toc}{chapter}{Resumen Ejecutivo}

EcoBrief Bolivia es una plataforma que utiliza inteligencia artificial de forma responsable para reducir el desperdicio digital informativo. El sistema recolecta noticias de medios bolivianos, elimina duplicados, prioriza el contenido relevante y genera res\'umenes claros para que las personas no tengan que abrir m\'ultiples p\'aginas, leer la misma historia varias veces o consumir datos innecesarios.

La propuesta no busca crear m\'as contenido. Busca reducir el ruido digital. EcoBrief transforma un flujo disperso de noticias, redes sociales y fuentes digitales en informaci\'on breve, priorizada y medible:

\begin{itemize}[nosep]
    \item De muchas p\'aginas a pocos res\'umenes.
    \item De noticias repetidas a historias \'unicas.
    \item De \textit{scroll} social a informaci\'on trazable.
    \item De IA aplicada sin filtro a IA aplicada despu\'es de deduplicar.
    \item De navegaci\'on manual a informaci\'on personalizada.
    \item De impacto ambiguo a m\'etricas visibles.
\end{itemize}

El prototipo ya cuenta con \textit{web scraping} real de medios locales, clasificaci\'on por categor\'ias, ranking de prioridad, res\'umenes con IA, deduplicaci\'on hist\'orica, base de datos PostgreSQL, frontend web y p\'agina de impacto. El sistema est\'a desplegado en un VPS de Hostinger y accesible p\'ublicamente.

\begin{center}
    \fcolorbox{verde-ecobrief}{verde-ecobrief!5!white}{\parbox{0.85\textwidth}{
        \centering\textbf{Mensaje central:}\\[0.3cm]
        \textit{EcoBrief Bolivia usa IA no para producir m\'as ruido, sino para reducirlo.}
    }}
\end{center}
